% =============================================================
% Revised quantitative section — addresses Sewon's two comments
% by re-grouping the relation table by audit role and by adding
% definitions / takeaways. Multi-hop dependency is left to the
% qualitative paragraphs and graph-query examples (intentionally:
% relation groups summarize single-edge roles; multi-hop model-
% data-model dependencies are surfaced by graph/path queries).
% =============================================================

\section{Analysis of Recovered Dependency Graphs}\label{sec:findings}

In this section, we analyze the graph recovered by \system{} and highlight findings that would be difficult to surface from individual sources alone. One advantage of representing recursive dependencies as a structured graph is that audit questions become executable queries rather than manual document searches. We can issue graph/SQL-style queries over the recovered graph: find released models that transitively depend on generators with restrictive licenses; identify benchmarks that appear both as evaluation targets and as upstream training seeds, prompts, or concept sources; locate operations supported by only one source type; or detect stages where papers, cards, and code provide conflicting evidence. The findings below were surfaced through this query-driven audit workflow: structured queries identify candidate risk patterns, and the attached source anchors make those candidates verifiable.

Across the analyzed releases, \system{} recovers a large, heterogeneous, and deeply recursive dependency graph. Table~\ref{tab:graph-scale} reports the scale and depth of the recovered lineage for each target model, while Table~\ref{tab:graph-composition} summarizes the aggregate graph composition. The graph contains both model and dataset artifacts, reflecting that model dependencies often enter indirectly through generated, filtered, or rewritten datasets rather than only through base-model inheritance.

\begin{table}[t]
\centering
\small
\caption{Scale and depth of recovered dependency graphs. Ancestors count unique transitive upstream artifacts reachable from each target model.}
\label{tab:graph-scale}
\begin{tabular}{lcc}
\toprule
Target model & Ancestors & Max depth \\
\midrule
Olmo-3-Instruct & 618 & 8 \\
Olmo-3-Think & 675 & 11 \\
Olmo-3-Base & 601 & 11 \\
Nemotron-3-Super & 667 & 6 \\
Nemotron-3-Nano-Base & 362 & 5 \\
DR-Tulu & 33 & 4 \\
SmolLM3-Base & 76 & 2 \\
\midrule
\textbf{Aggregate graph} & \multicolumn{2}{c}{2,844 nodes, 14,769 edges, 51,456 anchors} \\
\bottomrule
\end{tabular}
\end{table}

\noindent
The wide variance in max depth across targets reflects pipeline complexity rather than recovery quality: SmolLM3-Base is a pretraining-only release whose recoverable lineage is dominated by web-scrape and filter steps (depth 2), whereas the Olmo-3 Instruct/Think families pass through midtraining, multiple SFT/DPO/RL stages, and judge-mediated synthetic-data construction (depth 8--11). Lineage size is also lower-bounded by what upstream documentation makes recoverable; releases with shallower public artifact trees (e.g., DR-Tulu) yield correspondingly smaller ancestor sets.

\begin{table}[t]
\centering
\small
\caption{Aggregate composition of the recovered graph after deduplication.}
\label{tab:graph-composition}
\begin{tabular}{lr}
\toprule
Statistic & Count \\
\midrule
Source lattices merged & 193 \\
Relation artifacts merged & 715 \\
Artifact identity clusters & 2,545 \\
Artifact nodes & 2,844 \\
\quad Dataset nodes & 1,472 \; (51.8\%) \\
\quad Model nodes & 1,070 \; (37.6\%) \\
\quad Off-lattice nodes & 299 \; (10.5\%) \\
\quad Entity nodes & 3 \\
\midrule
Dependency edges & 14,769 \\
Evidence anchors & 51,456 \\
Flagged cross-source conflicts & 60,708 \\
\bottomrule
\end{tabular}
\end{table}

\noindent
Off-lattice nodes are low-degree edge endpoints retained for auditability but not promoted into canonical identity lattices. Flagged cross-source conflicts are candidate disagreements in identity attributes, descriptions, or anchor metadata; most are minor normalization conflicts (e.g., \texttt{size=7B} vs \texttt{7-billion}), while a smaller subset correspond to substantive documentation discrepancies analyzed below.

Table~\ref{tab:relation-roles} groups the recovered relations by the audit role each edge plays. Rather than reading off a flat per-relation count, this grouping organizes edges by what they support: training and evaluation inputs, upstream model operations on data (a model generating, filtering, transforming, or embedding a piece of data), weight-level model lineage, and methodology or audit-method relationships such as recipe inspiration and explicit decontamination. The grouping summarizes single-edge roles; multi-hop model--data--model dependencies (e.g., where a teacher model generated data that another model trained on) are not collapsed into a single relation type and are instead surfaced by the graph/path queries used in the qualitative findings below. This is supported by recent evidence that model-generated data can transmit model-specific behavioral traits even when explicit references are filtered out, including through code and reasoning traces~\cite{cloud2025subliminallearninglanguagemodels}.

\begin{table}[t]
\centering
\small
\caption{Recovered dependency edges grouped by audit role. Groups describe the role of an edge rather than the full transitive lineage; multi-hop model--data--model dependencies are surfaced by graph/path queries in the qualitative findings. Percentages are computed over all 14,769 recovered edges.}
\label{tab:relation-roles}
\begin{tabular}{p{0.30\linewidth}p{0.43\linewidth}rr}
\toprule
Audit role & Constituent relations & Count & Share \\
\midrule
Training / evaluation inputs
  & \texttt{used\_for\_evaluation}, \texttt{trained\_on}, \texttt{used\_for\_ablation}
  & 11,517 & 78.0\% \\
\addlinespace
Upstream model operations on data
  & \texttt{generated\_by}, \texttt{filtered\_by}, \texttt{transformed\_by}, \texttt{embedded\_by}
  & 1,996 & 13.5\% \\
\addlinespace
Weight-level model lineage
  & \texttt{trained\_from}, \texttt{merged\_from}, \texttt{quantized\_from}
  & 560 & 3.8\% \\
\addlinespace
Methodology / audit influence
  & \texttt{inspired\_by}, \texttt{decontaminated\_against}
  & 696 & 4.7\% \\
\midrule
\textbf{Total} & & \textbf{14,769} & 100\% \\
\bottomrule
\end{tabular}
\end{table}

\noindent
The four-group view makes one structural observation immediate: edges where upstream models generate, filter, transform, or embed data account for $13.5\%$ of recovered edges, compared with $3.8\%$ for direct weight-level model lineage. Thus, upstream models more often enter the graph by shaping data than by serving as base checkpoints, which is precisely the regime where multi-document, multi-hop tracing yields audit value.

The recovered evidence is also highly fragmented across source types. As shown in Table~\ref{tab:source-type-distribution}, most operations are supported by only one source class, meaning that systems restricted to a single document type would miss a large fraction of recoverable lineage.

\begin{table}[t]
\centering
\small
\caption{Source-type distribution of recovered operations. Single-source operations are supported by evidence from only one source class; multi-source operations are supported by anchors from multiple source classes.}
\label{tab:source-type-distribution}
\begin{tabular}{lrr}
\toprule
Source support & Count & Share \\
\midrule
Only Hugging Face cards & 885 & 6.0\% \\
Only training code / scripts & 3,712 & 25.2\% \\
Only PDFs / technical reports & 4,817 & 32.8\% \\
Only release blogs / docs & 61 & 0.4\% \\
Only other docs & 983 & 6.7\% \\
Multiple source types & 4,243 & 28.9\% \\
\midrule
Total operations & 14,701 & 100\% \\
\bottomrule
\end{tabular}
\end{table}

\noindent
After deduplication, $28.9\%$ of operations are supported by anchors spanning multiple source classes. The remaining $71.1\%$ are supported by a single source class, with PDFs / technical reports and training code accounting for the largest single-source shares. Because no source class covers a majority of operations, audits restricted to a single document type would miss many recoverable lineage facts and would also lose cross-source corroboration for multi-source operations.
